\section{Legal Framework}
\label{sec:legal}

\subsection{Reynolds v. Sims and Population Equality}

The constitutional foundation of redistricting law is \emph{Reynolds v. Sims}, 377 U.S. 533 (1964)~\citep{reynolds1964}, which held that the Equal Protection Clause requires legislative districts to contain approximately equal population. Chief Justice Warren's opinion did not specify which population---total persons, citizens, registered voters, or eligible voters---must be equalized. The Court held only that whatever population base is chosen, it must be applied equally across districts.

This deliberate ambiguity in \emph{Reynolds} has generated decades of litigation. In \emph{Burns v. Richardson}, 384 U.S. 73 (1966)~\citep{burns1966}, the Court addressed Hawaii's use of registered voters as its redistricting base to exclude large numbers of non-resident military personnel counted in the federal census. The Court approved Hawaii's approach, holding that the Equal Protection Clause does not compel any particular population base as long as the resulting plan does not ``systematically discriminate against any identifiable class.'' The \emph{Burns} decision is the principal precedent authorizing population-base alternatives to total population for \emph{state} redistricting.

\subsection{Evenwel v. Abbott (2016): The Controlling Authority}

\emph{Evenwel v. Abbott}, 578 U.S. 54 (2016), is the most recent Supreme Court word on population-base choice. The plaintiffs---voters in rural and suburban Texas senate districts---argued that using total population to equalize Texas senate districts violated the Equal Protection Clause because it diluted their votes: urban districts with large noncitizen populations had fewer eligible voters per district than rural districts with few noncitizens, yet each district carried the same representative.

The Court, in an opinion by Justice Ginsburg joined by seven justices, unanimously rejected this argument. The opinion rested on three pillars:

\begin{enumerate}
    \item \textbf{Historical practice}: The United States has counted total population---including noncitizens---for apportionment and redistricting purposes since the founding. The Fourteenth Amendment's Apportionment Clause (ratified 1868) counts ``the whole number of persons in each State,'' which has always included noncitizens.

    \item \textbf{Representation theory}: The Court endorsed the view that representatives represent all constituents---citizens and noncitizens alike---in their districts. Noncitizens use government services, are subject to laws, and have cognizable interests in the quality of local government, schools, and public safety. Equalizing total population ensures that each representative serves an approximately equal number of constituents in this broader sense.

    \item \textbf{No constitutional mandate for CVAP}: The Court explicitly refused to hold that citizen VAP is constitutionally required or that total population is constitutionally mandated. The opinion reserved the question whether states \emph{may} use citizen VAP instead of total population.
\end{enumerate}

Justice Thomas concurred in the judgment only, arguing that \emph{Reynolds} itself was wrongly decided and should be limited. Justice Alito concurred separately, suggesting sympathy for a future CVAP-permissibility holding. The two concurrences are significant: they signal that at least two justices---possibly more---would uphold a state's choice to use citizen VAP for redistricting if such a case reached the Court.

\subsection{Department of Commerce v. New York (2019)}

\emph{Department of Commerce v. New York}, 588 U.S. 752 (2019)~\citep{deptcommerce2019}, does not directly address redistricting population bases, but it is essential context. The Court blocked the Trump administration's addition of a citizenship question to the 2020 Census questionnaire, finding that the stated justification (VRA enforcement) was pretextual. The Court did not hold that citizenship data cannot be collected; it held that the \emph{manner} in which the decision was made---with a pretextual rationale---was arbitrary and capricious under the Administrative Procedure Act.

The \emph{Department of Commerce} litigation had a redistricting subtext that was widely understood: a citizenship question on the Census could be used to develop CVAP redistricting data at the census-block level, enabling states to draw districts on a citizen-VAP basis. The decision's blocking of the citizenship question did not close the door to CVAP redistricting---the Census CVAP special tabulation (derived from ACS data) remains available---but it removed the highest-resolution potential data source for such a plan.

\subsection{The Open Question After Evenwel}

\emph{Evenwel} resolved that states \emph{may} use total population. It left open:
\begin{enumerate}
    \item Whether states \emph{may} use citizen VAP instead
    \item Whether states \emph{may} use other bases (registered voters, eligible voters)
    \item Whether the federal government \emph{must} use total population for congressional apportionment (a distinct question addressed separately below)
\end{enumerate}

The first open question is the most immediately litigable. Several states have considered or enacted CVAP-based redistricting for state legislative maps; federal courts have reviewed some of these plans. No Supreme Court decision has resolved whether a state's election of citizen VAP over total population violates the Equal Protection Clause by systematically disadvantaging urban-immigrant communities.
